% PROVED/EXACT RESTRICTED; independent referee pending. Neither P=NP nor P!=NP proved.
\documentclass[10pt]{article}
\usepackage[margin=.72in]{geometry}
\usepackage{amsmath,amssymb,amsthm,booktabs,array,xurl,hyperref,seqsplit}
\setlength{\emergencystretch}{3em}
\newtheorem{theorem}{Theorem}
\newtheorem{lemma}[theorem]{Lemma}
\newtheorem{corollary}[theorem]{Corollary}
\newtheorem{proposition}[theorem]{Proposition}
\newcommand{\F}{\mathbb F_2}
\title{VARIABLE-LENGTH-SATURATION238:\\An Orthogonal-Array Obstruction to One-Bit Reusable Parity Summaries}
\author{LSPV-2, Circuit and Reduction Builder}
\date{29 October 2026, 20:17 America/Los\_Angeles}
\begin{document}
\maketitle
\begin{center}\fbox{\parbox{.94\linewidth}{\textbf{PROVED/EXACT, RESTRICTED, independent review pending.} In the precisely declared retained-source, two-stage indexed-bit-probe model, the optimal \emph{maximum description length} of a non-prefix, variable-length deterministic reusable summary of a six-bit source is exactly \(2\) bits when every future parity query must be answered with at most two adaptive raw-source bit probes. The lower bound allows unrestricted nonlinear preprocessing, unrestricted continuation computation, and three distinct length-at-most-one source-dependent summary strings. The proof strengthens the preceding degree-volume fiber bound by an orthogonal-array saturation obstruction. It is \emph{not} an unrestricted SAT lower bound. Both \(P=NP\) and \(P\ne NP\) remain \textbf{OPEN}.}}\end{center}

\section{PROPOSE: precise missing lemma and four falsifiers}
The preceding \texttt{TWO-PROBE-NONLINEAR-FRONTIER237} proved the tight fixed-length optimum \(L_2(6)=2\) but left the non-prefix, variable-length model open: can three messages \(\{\epsilon,0,1\}\) cover all 64 sources? The elementary fiber-volume bound gives \(V(6,2)=1+6+15=22\), and the inconclusive arithmetic \(3\cdot22=66\ge64\). We test four distinct mechanisms.
\begin{description}
\item[] \textbf{M1, unrestricted nonlinear partition:} split the 64 sources into three admissible fibers, each supporting every parity query by an adaptive two-probe decoder. \textbf{FALSIFIED/RESTRICTED} below by a structural upper bound of 21 points per fiber.
\item[] \textbf{M2, saturation-to-design transfer:} if a fiber achieves exactly \(V(b,Q)\) points, every parity's Boolean low-degree extension forces the degree-\(Q\) Walsh evaluation matrix to have mutually orthogonal columns. \textbf{PROVED/EXACT}.
\item[] \textbf{M3, maximal-length two-bit control:} a concrete affine cached representation uses exactly two persisted bits, reads at most one bit in each of two triples, and gives every parity. \textbf{PROVED/EXACT}; tested on all \(64^2\) source/query pairs.
\item[] \textbf{M4, SAT transport:} the storage bound excludes a joint SAT solver that reads its entire input, or a summary whose future context was known during preparation. \textbf{FALSIFIED/RETIRED} as an inference by direct input-scan and query-specific controls.
\end{description}

\section{SEARCH: project provenance and literature guard}
The actual project files inspected are \path{research/agents/lspv-2/2026-10-29_1912_LSPV-2_TWO_PROBE_NONLINEAR_FRONTIER237.tex}, blob \seqsplit{cba318698fad1541f0d9d187d2b607d12d44b61d}, and its executable, plus \path{research/agents/lspv-2/2026-10-29_1838_LSPV-2_NONLINEAR_ONE_PROBE236.tex} (prior blob \seqsplit{2d3deb8ee9496970dd12139860f88993efef5345}). \path{AGENTS.md} (\seqsplit{19656af102604d5f0a721d548cf2e4e72a063561}), \path{CODEX.md} (\seqsplit{754999f50de93cdb693ed37bc8110d31e9e698ef}) and \path{.codex/skills/pnp-research/SKILL.md} (\seqsplit{b8df83a7b0b1428f349b0e02ce1387d979e8db7f}) were checked. Root \texttt{main.tex} fetched as blob \seqsplit{df263c212ae0c46963c82e6f994537f0e81b590a}, with substantial text and an information-value section, but any root update needs a \emph{fresh} ownership-marker check.

The underlying polynomial representation of decision trees is classical~\cite{nisan1994}; the compiled-data representation-versus-query distinction is standard~\cite{darwiche2002}. The retained-original-plus-sketch model is particularly close to the systematic query-with-sketch literature; the September 2026 preprint of Garg et al.~\cite{garg2026} concerns a different matrix-powering problem, not the present exact parity fiber. Orthogonal arrays and Walsh-vanishing constraints are classical combinatorics~\cite{bierbrauer1995}. Their use in the following restricted equality test is independently proved, without asserting novelty over the broader coding-theoretic literature. Repositories searched: \texttt{j15marti/pnp} (named Q=2 theorem, exact prior executable, governance), \texttt{sagemath/sage} (coding directory for Hamming and Reed--Muller implementations); no Sage code was used as a theorem oracle.

\section{Declared two-stage model and source authentication}
Let \(b\ge1\). A source \(a=(a_1,\ldots,a_b)\in\F^b\) is prepared \emph{before} a future arbitrary query mask \(w\in\F^b\) is revealed. A deterministic preparer, possibly of unbounded time and computational complexity, outputs one source-dependent \emph{binary string} \(s(a)\) of length at most \(L\). The string length itself is visible, and the strings need not form a prefix-free code. No other source-dependent metadata, counters, transcript, output timing, or advice is persisted for free. The immutable full source \(a\) is still available to the continuation, but it may make at most \(Q\) indexed \emph{raw source-data bit} reads, adaptively. All querying of \(s(a)\) is permitted; continuation computation is unlimited for the lower bound. The decoder must return the \emph{exact} bit \(w\cdot a\bmod 2\) for every pair \((a,w)\).

One SAT realization is \(A_a(X)=\bigwedge_i (x_i=a_i)\) and \(B_{w,t}(X)=[w\cdot X=t]\). Then \(\mathrm{SAT}(A_a\wedge B_{w,t})\iff w\cdot a=t\). This is a unary-CNF-plus-affine-context model; it is neither a connected strict exact-$(2,2)$ grammar nor a proof that arbitrary SAT queries induce this two-stage access restriction. The source and query encodings, all indexes and their physical bit lengths, and all observed outputs have separate native-model charges.

\section{Saturation forces an orthogonal array}
Write \(V=V(b,Q)=\sum_{j=0}^{Q}\binom bj\), \(\chi_u(a)=(-1)^{u\cdot a}\), and \(\mathcal U=\{u\in\F^b:|u|\le Q\}\).
\begin{lemma}[Boolean saturation orthogonality]\label{lem:saturation}
Suppose \(S\subseteq\F^b\) has \(|S|=V(b,Q)\), and for every mask \(w\in\F^b\) there exists a Boolean function \(g_w:\F^b\to\{-1,1\}\) of real Walsh--Fourier degree at most \(Q\) such that \(g_w(a)=\chi_w(a)\) for every \(a\in S\). Then the column vectors \((\chi_u(a))_{a\in S}\), \(u\in\mathcal U\), are pairwise orthogonal. Consequently \(S\) is a binary orthogonal array of strength \(t=\min(2Q,b)\), and \(2^t\) divides \(|S|=V(b,Q)\).
\end{lemma}
\begin{proof}
Form the square \(V\times V\) evaluation matrix \(E=(\chi_u(a))_{a\in S,u\in\mathcal U}\), and the \(V\times2^b\) full-character matrix \(H=(\chi_w(a))_{a\in S,w\in\F^b}\). Every column of \(H\) is the restriction of a function with degree at most \(Q\), so it lies in the column space of \(E\). The full Walsh characters span all real functions on the finite set \(S\), hence \(\operatorname{rank} E=|S|=V\), and \(E\) is invertible. Writing \(c_w=E^{-1}H_{\cdot,w}\), the unique Fourier coefficient vector of \(g_w\) in degrees at most \(Q\) is \(c_w\). Booleanity of \(g_w\) on the \emph{entire cube}, not just \(S\), implies Parseval's identity \(\|c_w\|_2^2=1\).

Distinct rows of the full Walsh matrix are orthogonal, so \(HH^{\mathsf T}=2^b I_V\). Averaging Parseval over all query masks gives
\[
1=2^{-b}\sum_w\|E^{-1}H_{\cdot,w}\|_2^2
  =\operatorname{tr}\bigl((E^{\mathsf T}E)^{-1}\bigr).
\]
Let \(G=E^{\mathsf T}E\) with positive eigenvalues \(\lambda_i\). Each column of \(E\) has squared length \(V\), so \(\operatorname{tr}G=V^2\). By Cauchy--Schwarz,
\[
\operatorname{tr}(G^{-1})=\sum_i\lambda_i^{-1}
\ge\frac{V^2}{\sum_i\lambda_i}=1.
\]
Equality holds, hence all eigenvalues are \(V\) and \(G=VI_V\). For every nonempty \(T\subseteq[b]\) with \(|T|\le2Q\), select distinct \(u,v\in\mathcal U\) whose symmetric difference is \(T\). Orthogonality gives \(\sum_{a\in S}\chi_T(a)=0\). Fourier inversion on every set of \(t=\min(2Q,b)\) coordinates shows that every one of its \(2^t\) bit patterns occurs exactly \(|S|/2^t\) times. Integrality implies \(2^t\mid|S|\), proving the lemma.
\end{proof}

\begin{corollary}[non-prefix refined finite-message bound]\label{cor:cap}
In the two-stage deterministic source-probe model, every summary fiber has size at most
\[
 C(b,Q)=V(b,Q)-\mathbf 1\!\left[2^{\min(2Q,b)}\not\mid V(b,Q)\right],
\]
and every non-prefix, variable-length summary of maximum length \(L\) satisfies
\[
\boxed{(2^{L+1}-1)C(b,Q)\ge2^b.}
\]
These are necessary conditions, not claims that the fiber cap is always attained.
\end{corollary}
\begin{proof}
An adaptive depth-\(Q\) deterministic bit decision tree has a globally Boolean real multilinear representation of degree at most \(Q\), by induction through its queried variables. Therefore Lemma~\ref{lem:saturation} applies if any fiber reaches the preceding volume ceiling \(V\). If the stated divisibility condition fails, the ceiling is unattainable and every fiber has size at most \(V-1\). Otherwise the original volume ceiling remains. Finally there are exactly \(1+2+\cdots+2^L=2^{L+1}-1\) binary strings of length at most \(L\); the fibers partition all \(2^b\) source words.
\end{proof}

\section{Exact six-source-bit variable-length frontier}
\begin{theorem}[first non-prefix two-probe exact point]\label{th:6}
For \(b=6,Q=2\), if arbitrary deterministic nonlinear preprocessing may emit one non-prefix binary string of length at most \(L\), while every subsequent parity mask must be answered by a deterministic continuation making at most two adaptive indexed raw-source-bit probes, then the exact minimum \(L\) is \(2\).
\end{theorem}
\begin{proof}
The degree volume is \(V(6,2)=1+6+15=22\). If a fiber had 22 points, Lemma~\ref{lem:saturation} would force it to be an orthogonal array of strength four, which is impossible since \(16\not\mid22\). Thus every fiber has at most 21 elements. A maximum-one-bit non-prefix summary has only three labels \(\epsilon,0,1\), covering at most \(3\cdot21=63<64=2^6\) sources. Hence \(L\ge2\).

For the matching algorithm, divide the six source coordinates into triples \((a_1,a_2,a_3)\) and \((a_4,a_5,a_6)\). Persist the two bits \(s_1=a_1\oplus a_2\oplus a_3\) and \(s_2=a_4\oplus a_5\oplus a_6\). In each triple, the three original coordinates correspond to nonzero label vectors \(01,10,11\) in \(\F^2\), with the relation \(a_3=s_1\oplus a_1\oplus a_2\) (and analogously for the second triple). For a query mask restricted to a triple, form the XOR of the selected coordinate labels \(c\in\F^2\). Every requested parity is its known stored-bit correction plus either zero (if \(c=0\)) or one appropriate original bit (if \(c\ne0\)). This uses at most one raw bit probe per triple, hence at most two in total. It answers every \(w\in\F^6\) exactly with two stored bits. Preprocessing and continuation are uniform and polynomial. Thus \(L=2\) suffices.
\end{proof}

\begin{table}[h]\centering\small
\begin{tabular}{@{}rrrrrr@{}}\toprule
$b$&$Q$&$V(b,Q)$&original $(2^{L+1}-1)V$&refined count&conclusion\\\midrule
6&2&22&$3\cdot22=66$&$3\cdot21=63$&$L=1$ impossible; $L=2$ attained\\
8&2&37&$7\cdot37=259$&$7\cdot36=252$&$L=2$ impossible\\
11&2&67&$31\cdot67=2077$&$31\cdot66=2046$&$L=4$ impossible\\
\bottomrule
\end{tabular}
\caption{Exact integer arithmetic for the sharpened necessary bound. Only the first row includes an established matching upper construction; other rows do not assert optimality.}
\end{table}

\section{Charged construction--storage--continuation controls}
For \(b=6\), the native \emph{indexed bit-RAM} has a read-only encoded source \(A_a\) of bit length \(N=\Theta(b\log(b+1))\), separate read-only future query mask and output. The explicit constructor reads the six data bits, performs four pairwise XOR operations and writes two summary bits; general fixed-width source parsing and layout validation take \(O(N)\) time. The decoder scans each of the six query-mask bits, performs \(O(b)\) register XOR/address operations, reads at most two data bits, and may read stored bits. Generalizing the same two-block scheme to \(b=2(2^m-1)\), the established upper bounds are preprocessing \(O(N+bm)\) bit operations, per query \(O(N+bm)\), persisted source-dependent summary \(L=b-2m\) bits, and workspace \(O(L+\log N+m)\), with an immutable original source that remains \emph{physically present}. For the lower bound even unlimited time/peak workspace is allowed as long as the only source-dependent state carried across preparation is the declared summary.

In a two-tape read-only-input/writable-work implementation, an indexed source bit may require \(O(N)\) physical head movement, so the two indexed probes do not count as two native head steps. A one-read/write-tape implementation counts source cells, summary cells, carriage movement and rewrites in the same native footprint. Any fixed universal-machine interpreter must additionally charge its simulation overhead. The Python operation counters are direct interpreter events, not a hardware bit-flip or native TM crossing-sequence measurement. For \(M\) future queries, the explicit bit-RAM amortized upper ledger is \(O(N+bm+M(N+bm))\), including source parsing, all query scans, addressing, correction reads and answer output. There is no guessing, restart, hidden witness, or score-selecting subroutine in this specific constructor.

The source-free control retains all six original bits in the summary and uses zero subsequent source probes; a \emph{joint} solver that sees \(a,w\) simultaneously directly scans all six bits and uses no detached summary. A compiler knowing \(w\) beforehand can persist the one answer bit. These easy cases do not contradict Theorem~\ref{th:6} because they violate distinct, explicit stage/access constraints. No restriction is placed on a general SAT solver's workspace, representation or strategy.

\section{Exhaustive computation and contextual $(\mathrm{fact})\mid_v$ interpreter}
The dependency-free reproducer \texttt{2026-10-29\_2017\_LSPV-2\_VARIABLE\_LENGTH\_SATURATION238\_repro.py} checks all \(64\cdot64=4096\) pairs of six-bit source and parity mask against an independent direct parity evaluator. It also verifies the saturation divisibility arithmetic, enumerates a specific 22-point Hamming-ball control with full quadratic evaluation rank but nonorthogonal Gram columns, and enumerates SAT models on separate 2SAT, XOR, mixed-CNF, connected strict exact-$(2+,2-)$ three-variable SAT, and 15-variable UNSAT \(W_1\) controls. The last is exhaustively checked over all \(32768\) assignments. These finite tests support, but do not substitute for, Theorem~\ref{th:6}.

The minimal fact interpreter uses \((\mathrm{fact})\mid_v\) as a \emph{contextual score annotation}, not a conditional-probability operator. For a target parity \(T\), let \(\Phi_T(K)=\log_2|\{T(a):a\models K\}|\) for consistent certified affine facts \(K\). A new verified fact \(I\) receives \(v=\Phi_T(K)-\Phi_T(K\cup I)\), accompanied by explicit context, provenance, certificate, source-enumeration cost to discover the score, acquisition cost, verification cost and application cost. With target \(a_1\oplus a_2\), the facts \(a_1\oplus a_3=0\) and \(a_2\oplus a_3=0\) separately score zero and jointly score one; adding either again scores zero. This tests synergy and redundancy without charging a satisfying assignment as a free computational fact. Negative realized-entropy value is a separate scoring scale discussed in the previous report, not an additive SAT potential.

\section{Exact status, falsifiers, and next roles}
\noindent\textbf{PROVED/EXACT, RESTRICTED:} the general saturation-to-orthogonal-array lemma, refined necessary bound for non-prefix reusable summary messages, and exact \((b,Q,L)=(6,2,2)\) maximum-length frontier. \textbf{COMPUTATIONAL:} 4096 source/query comparisons, structural Gram and independent SAT/UNSAT tests, contextual scoring controls, Python same-run timing/memory counters. \textbf{KNOWN THEOREM:} Boolean decision-tree degree bound, Walsh orthogonality, basic orthogonal-array theory. \textbf{FALSIFIED/RETIRED:} three non-prefix messages of length at most one suffice for the six-bit, two-probe parity model. \textbf{OPEN:} exact non-prefix optimum for all other \(b,Q\), transfer to arbitrary connected strict22 SAT boundary functions, and any universal SAT time-space lower bound. No P-versus-NP resolution or algorithm-independent elapsed-time claim follows.

\noindent\textbf{LSPV-1:} independently referee the Parseval/equality argument and search for a saturated case where the OA divisibility is satisfied but a query decoder still fails. \textbf{LSPV-2:} seek a source-faithful, charge-preserving exact22 interface whose future contexts realize a sufficiently rich parity-query family, while allowing trivial affine solvers as controls. \textbf{LSPV-3:} examine simultaneous information-value scores for two correlated source facts and whether the saturation argument survives changing target contexts. \textbf{LSPV-4:} construct a native two-tape comparator for the six-bit scheme with measured head crossings, carriage moves and actual symbol changes. \textbf{LSPV-5:} referee the nonlinear decoder assumption, source-dependent side channels and root-main ownership; preserve both P versus NP universes as OPEN.

\begin{thebibliography}{9}
\bibitem{nisan1994} N. Nisan and M. Szegedy, ``On the Degree of Boolean Functions as Real Polynomials,'' \emph{Computational Complexity} 4 (1994), 301--313. doi:10.1007/BF01263419.
\bibitem{darwiche2002} A. Darwiche and P. Marquis, ``A Knowledge Compilation Map,'' \emph{JAIR} 17 (2002), 229--264. doi:10.1613/jair.989.
\bibitem{garg2026} S. Garg, S. He, Y. Li, P. A. Papakonstantinou, and X. Yang, ``Systematic Data Structure Lower Bounds via the Query-with-Sketch Model,'' arXiv:2609.18024 (September 2026).
\bibitem{bierbrauer1995} J. Bierbrauer, K. Gopalakrishnan, and D. R. Stinson, ``Orthogonal Arrays, Resilient Functions, Error-Correcting Codes, and Linear Programming Bounds,'' \emph{SIAM J. Discrete Math.} (1995). doi:10.1137/S0895480194270950.
\end{thebibliography}
\end{document}